% ============================================================
% Section 6 — Integration with the BX3 Framework
% ============================================================
\section{Integration with the BX3 Framework}
\label{sec:rs-integration}

The Recursive Spawning Protocol is not a module composed alongside the BX3 Framework; it is a \textit{structural instantiation} of the BX3 three-layer accountability architecture. Every RSP mechanism maps onto exactly one BX3 layer, and that mapping is not a design recommendation but a \textit{structural necessity}: the immutable parent pointer requires the Fact Layer's write-enforcement protocol, depth bounds require the Bounds Engine's constrained-execution evaluation, and the exclusive human terminus requires the Purpose Layer's non-negotiable authorization anchor. This section specifies the three mappings and establishes that no subset of layers achieves the same guarantees.

% ----------------------------------------------------------
\subsection{Purpose Layer: Accountability Anchor}
\label{sec:rs-purpose-anchor}

The Purpose Layer serves two structurally singular roles in RSP, neither of which any other layer can perform.

\medskip
\noindent\textbf{Exclusive origin of spawn authorization.} At root provisioning, the Purpose Layer defines the spawn authorization policy $P_{\mathrm{spawn}}$ and records it in the Fact Layer authorization ledger. $P_{\mathrm{spawn}}$ specifies two mandatory preconditions for any child spawn: (i) the child operating scope $R(n_{i+1})$ must be a strict subset of the parent scope $R(n_i)$ --- a descendant refinement, not a lateral expansion; and (ii) the spawn must be justified by operational necessity that cannot be met within the parent scope. This policy originates \textit{exclusively} at the Purpose Layer. The Bounds Engine evaluates proposed spawns against $P_{\mathrm{spawn}}$ but does not define it; the Fact Layer enforces it but does not originate it. No machine actor may issue a spawn warrant without a matching Purpose Layer authorization record.

\medskip
\noindent\textbf{Exclusive terminus of the delegation chain.} The chain terminus is defined as the human anchor $h(n_i)$ for all nodes $n_i$ in any chain $C$. Critically, $h(n_i)$ is non-collapsing: for all $i$, $h(n_i) = h(n_0)$. The human anchor does not migrate down the chain as depth increases. When the chain reaches $D_{\mathrm{ESCALATE}}$, the protocol compulsory-escalates to $h(n_0)$, who issues one of three directives --- \texttt{ACK}, \texttt{RESOLVE}, or \texttt{REJECT} --- and no machine actor may issue any of these. The chain terminates in human judgment, never in autonomous resolution. This is the structural expression of the upstream BX3 accountability guarantee: systems built on BX3 fail upward into human consciousness, never downward into autonomous opacity.

The structural necessity of the exclusive terminus is immediate. If any machine actor could absorb an exception or issue a terminal directive, the accountability chain would terminate at that actor rather than at a human. The Purpose Layer's exclusivity here is not a design preference; it is the load-bearing constraint that makes the upstream BX3 guarantee structurally operative rather than aspirational.

% ----------------------------------------------------------
\subsection{Bounds Engine: Depth and Scope Enforcement}
\label{sec:rs-bounds-depth}

The Bounds Engine provides constrained-execution evaluation: it evaluates every proposed action against the operating range defined by the Purpose Layer and enforces the depth and scope constraints that jointly bound the chain. The Bounds Engine cannot exceed these constraints under any operational circumstances.

\medskip
\noindent\textbf{Depth evaluation at spawn.} At each spawn event, the Bounds Engine evaluates the current chain depth $|C|$ against the Purpose-Layer-defined maximum $D_{\mathrm{max}}$, both anchored in the Fact Layer authorization ledger. If $|C| \geq D_{\mathrm{max}}$, the Bounds Engine does not execute the spawn. It transitions to \texttt{ESCALATING} state and initiates the Bailout Protocol, propagating the exception upward to $h(n_0)$. This enforcement is deterministic: the Bounds Engine has no discretion to exceed $D_{\mathrm{max}}$ regardless of urgency, operational exigency, or the importance of the unmet condition. The depth constraint is not a heuristic; it is a hard architectural boundary.

\medskip
\noindent\textbf{Scope constraint evaluation.} Every action $a$ proposed by a node must satisfy $a \in R(n)$, where $R(n)$ is the scope signature derived from the Purpose Layer at spawn time. The Bounds Engine performs this check before any action is dispatched to the Fact Layer for execution logging. An action that fails this check triggers a scope-violation exception, which is recorded in the Fact Layer and escalated if it persists.

\medskip
\noindent\textbf{Fact Layer pre-commit reinforcement.} The Bounds Engine's evaluations are structurally reinforced by the Fact Layer pre-commit hook. Before any spawn operation is recorded, the Fact Layer verifies two conditions: (1) the scope subset relationship $R(n_{i+1}) \subseteq R(n_i)$ holds, and (2) the resulting depth $\mathrm{depth}(n_{i+1}) \leq D_{\mathrm{max}}$ holds. The Fact Layer rejects any spawn operation violating either condition. The Bounds Engine cannot execute a rejected spawn; it may only transition to \texttt{ESCALATING}. The pre-commit hook converts depth bounds and scope refinement constraints from policy conventions into structural invariants --- overridable by no actor, under no circumstance.

The depth-limit insufficiency theorem (Section~\ref{sec:drift-depth-limits}) establishes that depth limits alone cannot prevent scope-creep drift. The Bounds Engine enforces the depth constraint; the Fact Layer enforces the scope refinement constraint; together they enforce both conjuncts of the drift-prevention condition. Neither layer alone is sufficient.

% ----------------------------------------------------------
\subsection{Fact Layer: Forensic Ledger}
\label{sec:rs-fact-ledger}

The Fact Layer provides deterministic enforcement and forensic auditability. It maintains the append-only ledger that makes the RSP chain auditable at any future time and enforces the write-protocol constraints that make RSP tamper-evident.

\medskip
\noindent\textbf{Append-only forensic ledger.} The Fact Layer records every spawn event, every state transition, every Purpose Layer directive received, and every unresolved exception. Each entry is timestamped, attributed to the node that generated it, and hash-chained to the preceding entry. Hash-chaining ensures that no entry can be inserted, deleted, reordered, or altered after recording: any tampering breaks hash continuity and is detectable by any observer with ledger read access. The ledger is the single source of truth for all authorized chain state --- it is the runtime artifact that distinguishes RSP chains from conventional delegation chains, which exist only as forensic reconstructions attempted after the fact.

\medskip
\noindent\textbf{Immutable parent pointer.} Once $\alpha(n)$ is established at node provisioning, no actor --- including the Purpose Layer after provisioning --- may modify it. The Fact Layer write protocol rejects any attempt to overwrite $\alpha(n)$. Chain topology changes occur only through authorized spawn operations recorded in the ledger; there is no mechanism for retroactive pointer manipulation. The parent pointer is not merely access-controlled; it is structurally immutable by protocol definition.

\medskip
\noindent\textbf{Pre-commit hook.} The pre-commit hook is the enforcement mechanism that makes all other RSP constraints structural. At every spawn event, it verifies two conditions before the spawn is recorded: (a) $R(n_{i+1}) \subseteq R(n_i)$ and (b) $\mathrm{depth}(n_{i+1}) \leq D_{\mathrm{max}}$. Rejected spawns are logged; the Bounds Engine does not execute them; the protocol transitions to \texttt{ESCALATING} if the unmet condition persists. The pre-commit hook operates identically on all machine actors; there is no privileged actor capable of bypassing it.

The structural necessity of the Fact Layer is established by the immutable-parent-pointer requirement: without the Fact Layer write protocol, immutability is a promise rather than a constraint. A mutable parent pointer enables all drift failure modes (Section~\ref{sec:drift-failure-modes}): silent orphaning, scope-creep drift, re-parenting drift, and ghost authority. Without the pre-commit hook, the scope refinement constraint and depth bound are policy documentation --- enforceable until an actor with sufficient privilege decides they are inconvenient. The Fact Layer is the structural enforcement layer that makes policy effective.

% ----------------------------------------------------------
\subsection{Structural Minimality}
\label{sec:rs-minimality}

The three-layer mapping is minimal and sufficient. Removing any layer causes the corresponding guarantee to fail, as summarized in Table~\ref{tab:layer-minimality}.

\begin{table}[h]
\centering
\begin{tabular}{p{3.2cm} p{5.4cm}}
\toprule
\textbf{Removed layer} & \textbf{Guarantee that fails} \\
\midrule
Purpose Layer          & Exclusive human terminus eliminated; chain can terminate in machine actors; spawn policy becomes self-authorized by machine actors \\
Bounds Engine          & Constrained execution evaluation eliminated; no layer evaluates operational necessity between pre-commit's binary approve/reject and mandatory escalation \\
Fact Layer             & Immutable parent pointer eliminated; retroactive chain manipulation possible; depth bound and scope refinement constraint reduce to policy conventions \\
\bottomrule
\end{tabular}
\caption{Layer removal consequences for RSP correctness properties.}
\label{tab:layer-minimality}
\end{table}

The three-layer architecture is the minimal structure that makes RSP's correctness properties unconditionally guaranteed rather than conditionally expected. No subset of layers achieves the same guarantees.
